\section{Simultaneous block replacement in finite groups}
\label{app:group-blocks}

\paragraph{Semigroup convention.}
The finite group $A$ and the morphisms into it are auxiliary to the shortening
proofs.  The semantic typing is always the positive-image restriction of the
window--group (or local-group) product, with the fresh identity used only to
evaluate empty template factors; by Lemma~\ref{lem:semigroup-transfer} this
changes neither nonempty tuple comparisons nor the selected witnesses.

This appendix proves the purely group-theoretic statement used in
Theorem~\ref{thm:group-typed-shortest} and Lemma~\ref{lem:local-group-shortest}.
The commutative case is a zero-sum argument; the general case needs an
order-homogenization argument because, for $f\ge2$, blocks of different
indices can interleave differently on different lines.

\begin{definition}[Block-replacement family]
\label{def:block-family}
Let $A$ be a finite group and $f\ge1$.  A \emph{block-replacement family} of
width $f$ over $A$ with index set $[m]$ consists of $d\le f$ \emph{lines}; on
line $c$, a left-to-right sequence of blocks $\beta_{c,1},\ldots,\beta_{c,q_c}$
and constants $a_{c,0},\ldots,a_{c,q_c}\in A$; for every block $\beta$ an
\emph{owner} $o(\beta)\in[m]$ and two values $x_\beta,x'_\beta\in A$; every
index owns at most $f$ blocks.  For $J\subseteq[m]$ put
\[
 \mathrm{val}_c(J):=a_{c,0}\,y_{c,1}\,a_{c,1}\cdots y_{c,q_c}\,a_{c,q_c},
 \qquad
 y_{c,j}:=
 \begin{cases}
 x'_{\beta_{c,j}},& o(\beta_{c,j})\in J,\\
 x_{\beta_{c,j}},& o(\beta_{c,j})\notin J.
 \end{cases}
\]
\end{definition}

\begin{lemma}[Block replacement]
\label{lem:block-replacement}
In a block-replacement family of width $f$ over $A$ there is a nonempty
$J\subseteq[m]$ with $\mathrm{val}_c(J)=\mathrm{val}_c(\emptyset)$ for every
line $c$, provided that either
\textup{(a)} $A$ is abelian and $m>\mathsf d(A^f)$, or
\textup{(b)} $A$ is arbitrary and $m\ge M_{A,f}$, where $\mathsf d$ and
$M_{A,f}$ are as in Section~\ref{subsec:group-invariants}.
\end{lemma}

Part (a) is proved in Appendix~\ref{app:window-abelian} and part (b) in
Appendix~\ref{app:window-group}.

\subsection{Commutative groups: zero sums}
\label{app:window-abelian}

\begin{proof}[Proof of Lemma~\ref{lem:block-replacement}(a)]
Write $A$ additively.  By commutativity
$\mathrm{val}_c(J)-\mathrm{val}_c(\emptyset)=\sum_{i\in J}\delta_{i,c}$, where
$\delta_{i,c}$ is the sum of $x'_\beta-x_\beta$ over the blocks $\beta$ of
line $c$ owned by $i$.  Put $\delta_i=(\delta_{i,c})_c\in A^d$.  A sequence
over $A^d$ of length $m>\mathsf d(A^f)\ge\mathsf d(A^d)$ has a nonempty
zero-sum subsequence $(\delta_i)_{i\in J}$; the inequality holds because $A^d$
embeds in $A^f$ and a zero-sum-free sequence in a subgroup is zero-sum free in
the group.
\end{proof}

\begin{remark}[Size of the zero-sum constant]
\label{rem:davenport-size}
The prefix-sum pigeonhole argument gives $\mathsf d(A^f)\le|A|^f-1$; the
zero-sum bound is usually much stronger.  Write
$A\cong C_{n_1}\oplus\cdots\oplus C_{n_s}$ and $n=\exp(A)\ge2$.  Then
\[
 f\sum_{i=1}^{s}(n_i-1)\ \le\ \mathsf d(A^f)\ \le\
 (n-1)+n\log\frac{|A|^f}{n},
\]
where the lower bound is attained by a basis sequence and the upper bound is
due to van Emde Boas and Kruyswijk
\cite[Proposition~5.1.8 and Theorem~5.5.5]{GeroldingerHalterKoch2006}; for a
$p$-group the lower bound is an equality
\cite[Theorem~5.5.9]{GeroldingerHalterKoch2006} (Olson).  Hence
$\mathsf d(A^f)=O(\exp(A)\,f\log|A|)$ is linear in $f$, whereas $|A|^f-1$ is
exponential; for $A=C_2$ the antichain bound drops from $2^f-1$ to $f$.  The
value $\mathsf d(A^d)$ is also the limit of the argument as stated, since a
zero-sum-free sequence of that length exists; we do not claim that a grammar
realizes such a configuration.
\end{remark}

\subsection{Arbitrary finite groups: order homogenization}
\label{app:window-group}

\begin{lemma}[A homogeneous block-order lemma]
\label{lem:group-block-order}
Fix positive integers $e,n$ with $n\ge3$.  For each
$t\in[e]$ and $i\in[n]$, let $(t,i)$ be a distinct point
on a line, with a total left-to-right order.  Suppose that
(i) the order of $(t,i)$ and $(u,i)$ is independent of $i$;
and (ii) for each ordered pair $(t,u)$, including $t=u$,
the order of $(t,i)$ and $(u,j)$ for $i<j$ is independent
of the choice of $i<j$.
Then the types at the points have a partition into classes,
each occupying a consecutive region of the line.  Every
class has one direction (increasing or decreasing in $i$);
inside its region the points occur in $n$ successive
\emph{chunks}, one per rank in that direction.  Each
chunk consists of exactly one point of each type in the
class (in a fixed within-chunk order).
\end{lemma}
\begin{proof}
By (ii) with $t=u$, each type's points are monotone.  For distinct types
$t,u$ write $X$ for $(t,i)<(u,j)$ and $Y$ for $(u,i)<(t,j)$, $i<j$; both are
uniform by (ii).  Suppose first that both types increase.  $X\wedge Y$ puts
all points of rank $i$ before all points of rank $j>i$.  Under $X\wedge\neg Y$
all off-diagonal comparisons put $t$ left of $u$, and so do equal ranks:
$(u,i)<(t,i)$ would give $(t,1)<(t,2)<(u,1)<(t,1)$; thus the types are
separated.  The case $\neg X\wedge Y$ is symmetric, and $\neg X\wedge\neg Y$
gives $(t,1)<(t,2)<(u,1)<(u,2)<(t,1)$, a contradiction.  If $t$ increases and
$u$ decreases, $X\wedge Y$ gives $(t,2)<(u,3)<(u,1)<(t,2)$ and
$\neg X\wedge\neg Y$ gives $(t,3)<(u,2)<(t,1)<(t,3)$; under $X\wedge\neg Y$ all
off-diagonal comparisons put $t$ left of $u$, and a contrary equal-rank
comparison would give $(t,1)<(u,2)<(u,1)<(t,1)$, contrary to (i), so the types
are separated; $\neg X\wedge Y$ is symmetric.  Reversing indices covers two
decreasing types.  Thus any two types are either \emph{separated}, or have the
same direction and are \emph{interleaved by rank}.

Interleaving is transitive: if $t$ interleaves with $u$ and $u$ with $v$, all
three have one direction, say increasing; if $t,v$ were separated with $t$ to
the left we would have $(v,1)<(u,2)<(t,3)<(v,1)$, and with $v$ to the left
$(t,1)<(u,2)<(v,3)<(t,1)$.  Hence interleaving, with equality, partitions the
types into classes.  A type outside a class is separated from each type inside
it and lies on the same side of all of them, since otherwise two interleaved
types would be separated.  So class regions are consecutive, the rank
comparisons give successive chunks, and the within-chunk order is independent
of rank by (i).
\end{proof}

\begin{lemma}[Simultaneous finite-group interval shortening]
\label{lem:group-interval}
Let $A$ be any finite group, and consider a block-replacement family of width
$f$ over $A$ with $n\ge\max\{3,|A|^f\}$ indices in which all indices have the
same number of blocks, the same lines for them and the same within-index
order, and in which every cross-index comparison between the $t$-th block of
index $i$ and the $u$-th block of index $j$ on a common line has the same
truth value for all $i<j$.  Then a nonempty interval $J$ of indices satisfies
$\mathrm{val}_c(J)=\mathrm{val}_c(\emptyset)$ for every line $c$.
\end{lemma}
\begin{proof}
Apply Lemma~\ref{lem:group-block-order} on each line to the block types going
to that line.  Across all lines there are $q\le f$ class regions.  For such a
region $K$, enumerate its chunks $1,\ldots,n$ from left to right, including
the constants between the blocks of the region.  Let $P_p$ be the original
value of its prefix ending after chunk $p$ and $Q_p$ the value of that prefix
when chunks $1,\ldots,p$ are replaced; set $P_0=Q_0=1_A$ and
$\kappa^K_p=P_pQ_p^{-1}$.  Replacing exactly chunks $p+1,\ldots,p'$ gives the
ordered region value $P_pQ_p^{-1}Q_{p'}P_{p'}^{-1}P_n$, because
$Q_p^{-1}Q_{p'}$ is the replaced value of the segment between the two prefix
boundaries and $P_{p'}^{-1}P_n$ is the unchanged suffix.  Hence the region
keeps its value $P_n$ whenever $\kappa^K_p=\kappa^K_{p'}$; inverses are used,
but factors are never commuted.

If $K$ increases with the index order, the chunks of indices $(a,b]$ occupy
$(a,b]$; if it decreases, they occupy $(n-b,n-a]$.  For $j=0,\ldots,n$ define
a state in $A^q$ by taking $\kappa^K_j$ from each increasing and
$\kappa^K_{n-j}$ from each decreasing class.  There are $n+1>|A|^q$ states, so
two coincide, at $a<b$.  Replacing the indices $a+1,\ldots,b$ then preserves
every region value; everything else on every line is unchanged.
\end{proof}

\begin{proof}[Independent Ramsey proof of the qualitative assertion in Lemma~\ref{lem:block-replacement}(b)]
We first establish existence of a simultaneous replacement set without
using the quantitative order-homogenization bound below.  Let
$N=\max\{3,|A|^f\}$, and let $R(N,c)$ denote a finite two-uniform
Ramsey number for $c$ colors.  An owner with no blocks can be selected
alone.  Otherwise, among sufficiently many owners, select
$R(N,2^{f^2})$ of the same unary pattern (number of blocks, assigned
lines, and within-owner order); there are at most $U_f$ patterns.
Absorb all other owners into the unchanged constants.

The common pattern has $e\le f$ block types.  Color a pair of selected
owners $i<j$ by the truth values of
$p_t(i)<p_u(j)$ for all ordered pairs of types $(t,u)$ on the same
line.  There are at most $2^{e^2}\le2^{f^2}$ colors.  By finite Ramsey,
there is an ordered set of $N$ owners on which every comparison has a
fixed truth value, independently of $i<j$.  Within-owner comparisons
are already fixed by the common pattern.  On each line, the argument
of Lemma~\ref{lem:group-block-order} decomposes this homogeneous
configuration into consecutive rank-ordered regions, all increasing
or decreasing.  Lemma~\ref{lem:group-interval} now supplies a nonempty
interval of owners whose simultaneous replacement leaves the ordered
product on every line unchanged.

This proves the \emph{existence} assertion for an arbitrary finite,
possibly nonabelian group, without assuming any commutation of factors.
Ramsey's bound does not supply the asserted polynomial dependence on
$|A|$ for fixed $f$; that stronger, explicit assertion is established
next by the monotone-comparison extraction.
\end{proof}

\begin{lemma}[Polynomial order homogenization]
\label{lem:order-homogenize}
Let $N\ge3$, and define $m_0=N$ and
$m_{s+1}=(m_s-1)^2+1$, so that $m_k\le N^{2^k}$.
Let $T$ be a finite set of point types, each assigned to
one of several lines.  For every index $i\in[n]$ and
$t\in T$ let $p_t(i)$ be a point on the line of $t$, the
points on each line being distinct.  Let $\Pi$ be the set
of ordered pairs $(t,u)$, including $t=u$, of types on a
common line, and put $k=|\Pi|$.  If $n\ge m_k$, there is a
set $I\subseteq[n]$ of $N$ indices such that, for every
$(t,u)\in\Pi$, the truth value of $p_t(i)<p_u(j)$ is the
same for all $i<j$ in $I$.
\end{lemma}
\begin{proof}
Treat the pairs one at a time, the pairs $(t,t)$ first.  Each step passes
from a current index set of size at least $(m-1)^2+1$ to a subset of size at
least $m$; homogeneity already obtained, and monotonicity of each type,
persist in subsets.  For $(t,t)$, homogeneity means that $p_t$ is monotone, and
the Erd\H{o}s--Szekeres theorem gives a monotone subsequence of length $m$.

For $t\ne u$, both $p_t$ and $p_u$ are already monotone.  Relabel the current
indices $1,\ldots,m'$ in order, with $m'\ge(m-1)^2+1$, and put
$X(i,j)=[p_t(i)<p_u(j)]$ for $i<j$.  Each row $X(i,\cdot)$ and each column
$X(\cdot,j)$ is monotone; replacing $X$ by $1-X$ if necessary, the rows go
from $0$ to $1$.

If the columns go from $1$ to $0$, put $J(i)=\min\{j>i:X(i,j)=1\}$ (or
$\infty$); $J$ is nondecreasing.  Set $i_1=1$ and $i_{s+1}=J(i_s)$ while this
is at most $m'$.  For $a<b$, $i_b\ge J(i_a)$ and so $X(i_a,i_b)=1$; if the
chain has $m$ terms we are done.  Otherwise its at most $m-1$ terms cut $[m']$
into at most $m-1$ intervals $[i_s,i_{s+1})$ and $[i_L,m']$, one of which has
more than $m'/(m-1)>m-1$ indices; inside it every $i<j$ satisfies
$j<J(i_s)\le J(i)$, so $X(i,j)=0$ (and $J\equiv\infty$ in the last interval).

If the columns go from $0$ to $1$, the set $\{X=1\}$ is closed upward in both
coordinates.  Let $\kappa$ be the least index with $X(\kappa,\kappa+1)=1$, or
$\kappa=m'$.  For $\kappa\le i<j$, $X(i,j)\ge X(\kappa,\kappa+1)=1$; for
$i<j\le\kappa$, $X(i,j)\le X(j-1,j)=0$.  One of $[1,\kappa]$ and $[\kappa,m']$
has at least $(m'+1)/2\ge m$ indices, because $(m-1)^2+1\ge2m-1$.

After all $k$ pairs at least $m_0=N$ indices remain, and
$m_{s+1}\le m_s^2$ gives $m_k\le N^{2^k}$.
\end{proof}

\begin{proof}[Proof of Lemma~\ref{lem:block-replacement}(b)]
If an index owns no blocks, choosing that index alone preserves every
line, so assume each index owns at least one block. The
\emph{unary pattern} of an index records the number of its blocks, the
line of each, and their order within any shared line; there are at most
$f\cdot f^f\cdot f!=U_f$ patterns.  Hence at least $N_{A,f}^{\,2^{f^2}}$
indices share one pattern; the blocks of all other indices are never
replaced and are absorbed into the constants.  Number the blocks of the common
pattern as types $1,\ldots,f'$ ($f'\le f$) and use their positions as points
on their lines.  There are at most $f^2$ ordered pairs of types on a common
line, including $t=u$, and $m_k\le N^{2^k}\le N^{2^{f^2}}$ for $k\le f^2$, so
Lemma~\ref{lem:order-homogenize} with $N=N_{A,f}$ yields $N_{A,f}$ indices with
constant cross-index comparisons for all $i<j$.  These satisfy the hypotheses of
Lemma~\ref{lem:group-interval}, which supplies the required nonempty interval.
\end{proof}

For fixed $f$ the constant $M_{A,f}=U_f\max\{3,|A|^f\}^{2^{f^2}}$ is
polynomial in $|A|$; it replaces the multicolor Ramsey bound
$U_f\,c_f^{\,c_f(N_{A,f}-1)+1}$, $c_f=2^{f^2}$, of earlier versions, which is
exponential in $|A|^f$.  The dependence on $f$ remains doubly exponential.